\chapter{欠驱动系统与摆起控制：从部分反馈线性化到能量整形}
\label{ch:underactuated}

% ════════════════════════════════════════════════════════════════
%  控制部（归卷四《控制理论基础》）。
%  主源：Spong, Hutchinson, Vidyasagar《Robot Modeling and Control》2e，Ch.13。
%  集大成：以 Spong 为骨架，融入 Lyapunov/LaSalle（ch:lyapunov-stability）、
%          多输入反馈线性化（sec:ctrl-mimo-fbl）、LQR（sec:ctrl-lqr）、
%          计算力矩（sec:ctrl-robot-ctc）的统一视角。
%  右扰动/Hamilton 约定与全书一致；本章为欧氏状态空间 R^{2n} 上的机械系统控制。
% ════════════════════════════════════════════════════════════════

\begin{practice}[前置自测]
开始之前，先试答下列问题。若已能流畅作答，可只速读本章的正规型速查表与小结；若卡住，正是本章要为你打磨的。
\begin{enumerate}
  \item 全驱动机械臂可用计算力矩法（\cref{sec:ctrl-robot-ctc}）把非线性动力学整体抵消成 $\ddot{\mathbf{q}}=\mathbf{a}$。若执行器数少于自由度数，这套方法为何\emph{原封不动}用不了？最多还能线性化掉哪一部分？
  \item 一节单摆挂在小车上、小车有水平推力而摆轴无力矩，这个系统在\emph{倒立}平衡点附近线性化后可控吗？在\emph{下垂}平衡点附近呢？两者的差别从何而来？
  \item 倒立摆“甩起来（swing-up）”靠的是把能量一点点泵进系统。如何用一个 Lyapunov 函数把“注入能量”这件事变成一条可证明收敛的反馈律，而不是靠手调轨迹？
  \item 把欠驱动系统某个输出（如被动关节角）反馈线性化后，剩下的“内动态”可能是一个无阻尼单摆。这对你能否让该输出渐近镇定意味着什么？“非最小相位”在这里具体指什么？
  \item 为什么单靠能量法甩到倒立点附近还不够、必须在倒立点\emph{切换}到一个局部线性控制器（如 LQR）才能稳稳站住？切换的“捕获域”由谁决定？
\end{enumerate}
\end{practice}

\section*{本章目标}
学完本章，你应当能够：
\begin{enumerate}
  \item \textbf{识别并建模欠驱动系统}：写出欠驱动度 $\ell=n-m$，把动力学整理成上驱动/下驱动分块形式，认出小车倒摆、Acrobot、Pendubot、反应轮摆等典型配置；
  \item \textbf{在平衡点判线性可控性}：算出机械系统在平衡点的线性化 $(\mathbf{F},\mathbf{G})$，用 Kalman 秩判据判可控，并理解“被动关节须有非零势力”这一必要条件的物理含义；
  \item \textbf{做部分反馈线性化（PFL）}：推导配置型与非配置型 PFL，写出二阶正规型，掌握强惯量耦合条件；
  \item \textbf{分析零动态与最小相位性}：对输出 $\mathbf{y}=\mathbf{q}_2$ 或 $\mathbf{y}=\mathbf{q}_1$ 算出零动态，判断系统是否非最小相位；
  \item \textbf{设计能量整形 swing-up}：以总能量误差平方为 Lyapunov 候选，构造能量注入律把摆从下平衡泵向上平衡，并处理输入饱和；
  \item \textbf{搭切换控制}：用“远处能量法 $+$ 近处 LQR”的切换策略完成全局甩起加局部平衡，理解捕获域与切换面的角色。
\end{enumerate}

\section{从全驱动到欠驱动：问题的提出}\label{sec:ua-intro}

前面几章里，机械臂始终是\emph{全驱动}的：每个广义坐标都配一个独立执行器，动力学
\begin{equation}\label{eq:ua-fullact}
  \mathbf{M}(\mathbf{q})\ddot{\mathbf{q}}+\mathbf{C}(\mathbf{q},\dot{\mathbf{q}})\dot{\mathbf{q}}+\mathbf{g}(\mathbf{q})=\boldsymbol{\tau}
\end{equation}
中输入 $\boldsymbol{\tau}\in\mathbb{R}^n$ 与自由度数同维。这一“奢侈”带来了一系列强性质：系统\emph{全局可反馈线性化}，计算力矩法（\cref{sec:ctrl-robot-ctc}）令 $\boldsymbol{\tau}=\mathbf{M}(\mathbf{q})\mathbf{a}+\mathbf{C}\dot{\mathbf{q}}+\mathbf{g}$ 即把闭环化为 $n$ 个解耦的双积分器 $\ddot{\mathbf{q}}=\mathbf{a}$，任意可达轨迹都能跟踪。多输入反馈线性化（\cref{sec:ctrl-mimo-fbl}）从微分几何角度给出了同一结论的一般框架。

\begin{definition}[欠驱动机械系统]\label{def:ua-underactuated}
一个广义坐标 $\mathbf{q}\in\mathbb{R}^n$、独立控制输入 $\mathbf{u}\in\mathbb{R}^m$ 的拉格朗日机械系统，若 $m<n$，则称为\textbf{欠驱动}的。差 $\ell=n-m$ 称为系统的\textbf{欠驱动度}。被驱动的坐标称\emph{主动关节}，不被直接驱动的称\emph{被动关节}。
\end{definition}

欠驱动从多种渠道进入机器人学，理解其来源有助于把分散的现象归到同一框架下：
\begin{itemize}
  \item \textbf{有意设计}：故意去掉部分执行器以研究本质控制难题或降低成本。Acrobot（体操机器人）、Pendubot、反应轮摆都属此类。
  \item \textbf{模型本身引入}：把关节柔性纳入模型时，$n$ 连杆柔性关节机器人有 $2n$ 个自由度却只有 $n$ 个电机，天然欠驱动（这是上一章柔性关节模型的另一种解读）。
  \item \textbf{约束诱发}：单边接触约束（足底与地面）、非完整约束（轮子纯滚动）都使系统欠驱动。腿足机器人因脚地接触不可“施力拉地”而\emph{本质欠驱动}；轮式、游动、空间、飞行机器人同理。
  \item \textbf{动量守恒}：自由漂浮的空间机器人、空中翻转的体操/跳水机器人，其总动量守恒构成一条不可直接驱动的“坐标”。
\end{itemize}

\begin{insight}[一句话抓住难点]
全驱动系统“想去哪就去哪”，欠驱动系统“只能借自然动力学的力把自己甩过去”。前者的控制是\emph{抵消}动力学，后者的控制是\emph{利用}动力学——这正是为什么能量、无源性、零动态这些刻画“系统自身倾向”的工具在本章唱主角，而单纯的反馈线性化退居其次。
\end{insight}

本章以 Spong--Hutchinson--Vidyasagar 的处理为骨架，依次铺开三件工具：\textbf{部分反馈线性化}（把能线性化的那部分先线性化）、\textbf{能量/无源性方法}（驱动甩起）、\textbf{切换控制}（缝合全局甩起与局部平衡）。结尾将本章与另一类受限系统——非完整系统（\cref{ch:nonholonomic}）——以及安全控制的 CLF/CBF 框架（\cref{ch:clf-cbf}）接上。

\section{建模：上驱动与下驱动分块形式}\label{sec:ua-model}

一个一般的 $n$ 自由度欠驱动系统，动力学（推导工具同 \cref{ch:arm-dyn}）写成
\begin{equation}\label{eq:ua-dyn}
  \mathbf{M}(\mathbf{q})\ddot{\mathbf{q}}+\mathbf{C}(\mathbf{q},\dot{\mathbf{q}})\dot{\mathbf{q}}+\boldsymbol{\phi}(\mathbf{q})=\mathbf{B}\mathbf{u},
\end{equation}
其中 $\mathbf{M}(\mathbf{q})$ 是 $n\times n$ 对称正定惯量阵，$\boldsymbol{\phi}(\mathbf{q})$ 收集一切由势能导出的广义力（重力、弹性力等，即势能梯度），$\mathbf{B}$ 是 $n\times m$、秩为 $m$ 的输入分配矩阵，反映 $m$ 个独立执行器。为方便分析，把坐标重新编号，使主动关节聚在一起，则 $\mathbf{B}$ 取最简单的两种形式之一：
\begin{equation}\label{eq:ua-Bform}
  \mathbf{B}_u=\begin{bmatrix}\mathbf{I}_m\\\mathbf{0}\end{bmatrix}\ (\text{上驱动}),\qquad
  \mathbf{B}_l=\begin{bmatrix}\mathbf{0}\\\mathbf{I}_m\end{bmatrix}\ (\text{下驱动}).
\end{equation}

\begin{definition}[上驱动与下驱动系统]\label{def:ua-upperlower}
若前 $m$ 个（近端）关节主动、后 $\ell$ 个（远端）关节被动，称\textbf{上驱动系统}（取 $\mathbf{B}=\mathbf{B}_u$）；若后 $m$ 个关节主动、前 $\ell$ 个关节被动，称\textbf{下驱动系统}（取 $\mathbf{B}=\mathbf{B}_l$）。“上/下”取自“上臂/下臂”的类比而非关节编号。同一系统总可通过重排坐标在两种形式间互换，依分析方便择用。
\end{definition}

把 $\mathbf{q}=(\mathbf{q}_1,\mathbf{q}_2)$ 按被动/主动分块。以\textbf{下驱动}为例（$\mathbf{q}_1\in\mathbb{R}^\ell$ 被动、$\mathbf{q}_2\in\mathbb{R}^m$ 主动），\cref{eq:ua-dyn} 展开为
\begin{align}
  \mathbf{M}_{11}\ddot{\mathbf{q}}_1+\mathbf{M}_{12}\ddot{\mathbf{q}}_2+\mathbf{c}_1+\boldsymbol{\phi}_1&=\mathbf{0},\label{eq:ua-lower1}\\
  \mathbf{M}_{21}\ddot{\mathbf{q}}_1+\mathbf{M}_{22}\ddot{\mathbf{q}}_2+\mathbf{c}_2+\boldsymbol{\phi}_2&=\mathbf{u},\label{eq:ua-lower2}
\end{align}
其中惯量阵分块
\begin{equation}\label{eq:ua-Mblock}
  \mathbf{M}(\mathbf{q})=\begin{bmatrix}\mathbf{M}_{11}&\mathbf{M}_{12}\\\mathbf{M}_{21}&\mathbf{M}_{22}\end{bmatrix},\quad
  \mathbf{M}_{11}\in\mathbb{R}^{\ell\times\ell},\ \mathbf{M}_{22}\in\mathbb{R}^{m\times m},\ \mathbf{M}_{21}=\mathbf{M}_{12}^\top,
\end{equation}
$\mathbf{c}_i(\mathbf{q},\dot{\mathbf{q}})$ 为对应的科氏/离心项，$\boldsymbol{\phi}_i(\mathbf{q})$ 为势力项，$\mathbf{u}\in\mathbb{R}^m$ 是主动关节的输入广义力。上驱动形式（$\mathbf{q}_1\in\mathbb{R}^m$ 主动、$\mathbf{q}_2\in\mathbb{R}^\ell$ 被动）对称地把 \cref{eq:ua-lower1} 右端换成 $\mathbf{u}$、\cref{eq:ua-lower2} 右端换成 $\mathbf{0}$。

\begin{note}[被动行就是二阶约束]
下驱动的第一行 \cref{eq:ua-lower1} 右端恒为零。它不是“随便一条方程”，而是对运动施加的一条\textbf{二阶（加速度层）约束}：任何参考轨迹 $\mathbf{q}^d(t)$ 都必须满足它。于是欠驱动系统\emph{不能}跟踪任意轨迹——这是它与全驱动系统的又一根本区别。被动行扣住了 $\ddot{\mathbf{q}}_1$ 与 $\ddot{\mathbf{q}}_2$ 的耦合关系，正是后面 PFL 与零动态分析的入口。
\end{note}

\begin{example}[柔性关节机器人即欠驱动]\label{ex:ua-flexjoint}
$n$ 连杆柔性关节模型
\[
  \mathbf{D}(\mathbf{q}_1)\ddot{\mathbf{q}}_1+\mathbf{C}(\mathbf{q}_1,\dot{\mathbf{q}}_1)\dot{\mathbf{q}}_1+\mathbf{g}(\mathbf{q}_1)+\mathbf{K}(\mathbf{q}_1-\mathbf{q}_2)=\mathbf{0},\quad
  \mathbf{J}\ddot{\mathbf{q}}_2+\mathbf{K}(\mathbf{q}_2-\mathbf{q}_1)=\mathbf{u}
\]
恰是下驱动形式 \cref{eq:ua-lower1,eq:ua-lower2}，其中 $\mathbf{M}_{11}=\mathbf{D}$、$\mathbf{M}_{12}=\mathbf{M}_{21}=\mathbf{0}$、$\mathbf{M}_{22}=\mathbf{J}$、$\boldsymbol{\phi}_1=\mathbf{g}(\mathbf{q}_1)+\mathbf{K}(\mathbf{q}_1-\mathbf{q}_2)$。连杆角 $\mathbf{q}_1$（$2n$ 维状态里的被动部分）与电机角 $\mathbf{q}_2$ 只能指定其一的期望轨迹，另一个由动力学隐式决定。值得一提：柔性关节机器人是少数\emph{仍可全局反馈线性化}的欠驱动系统（取连杆角及其各阶导数为新坐标），这是个例外而非常态。
\end{example}

\section{典型欠驱动系统}\label{sec:ua-examples}

下面几个“最小模型”将贯穿全章，用来落地各概念。约定关节角用 DH 习惯，$\theta$（或 $q_1$）相对竖直方向度量。

\subsection{小车倒立摆（cart-pole）}\label{sec:ua-cartpole}
小车质量 $m_c$ 沿水平 $x$ 方向受输入力 $F$；长 $\ell$、末端质量 $m_p$ 的摆挂在小车上，摆轴\emph{无}力矩。它在运动学上等同一个 PR 两自由度机构。取 $x$ 为小车位置、$\theta$ 为摆相对竖直的偏角，欧拉--拉格朗日方程为
\begin{align}
  (m_p+m_c)\ddot{x}+m_p\ell\cos\theta\,\ddot{\theta}-m_p\ell\dot{\theta}^2\sin\theta&=F,\label{eq:ua-cart1}\\
  m_p\ell\cos\theta\,\ddot{x}+m_p\ell^2\ddot{\theta}-m_p g\ell\sin\theta&=0.\label{eq:ua-cart2}
\end{align}
取 $q_1=x,\,q_2=\theta,\,u=F$ 即\emph{上驱动}形式；若改取 $q_1=\theta,\,q_2=x$，同一系统又成\emph{下驱动}形式。小车倒摆是研究非线性动力学与平衡控制的经典原型：天车吊重的防摆、火箭竖直上升的俯仰、双足行走的平衡，都可抽象成它。

\subsection{Acrobot}\label{sec:ua-acrobot}
Acrobot（Acrobatic Robot）是一个\emph{仅第二关节有驱动}的两连杆 RR 机器人，模拟单杠上的体操运动员：$q_2,u$ 是髋角与髋力矩，握杠的第一关节无驱动。其动力学即两连杆 RR 机器人方程在第一关节力矩置零下的结果，是\emph{下驱动}形式：
\begin{align}
  m_{11}\ddot{q}_1+m_{12}\ddot{q}_2+c_1+g_1&=0,\label{eq:ua-acro1}\\
  m_{21}\ddot{q}_1+m_{22}\ddot{q}_2+c_2+g_2&=u,\label{eq:ua-acro2}
\end{align}
其中（$\ell_{ci}$ 为质心距、$I_i$ 为转动惯量）
\begin{equation}\label{eq:ua-acro-terms}
\begin{aligned}
  m_{11}&=m_1\ell_{c1}^2+m_2(\ell_1^2+\ell_{c2}^2+2\ell_1\ell_{c2}\cos q_2)+I_1+I_2,\\
  m_{12}&=m_{21}=m_2(\ell_{c2}^2+\ell_1\ell_{c2}\cos q_2)+I_2,\qquad m_{22}=m_2\ell_{c2}^2+I_2,\\
  c_1&=-m_2\ell_1\ell_{c2}\sin q_2\,(2\dot{q}_1\dot{q}_2+\dot{q}_2^2),\qquad c_2=m_2\ell_1\ell_{c2}\dot{q}_1^2\sin q_2,\\
  g_1&=(m_1\ell_{c1}+m_2\ell_1)g\cos q_1+m_2\ell_{c2}g\cos(q_1+q_2),\quad g_2=m_2\ell_{c2}g\cos(q_1+q_2).
\end{aligned}
\end{equation}

\subsection{Pendubot}\label{sec:ua-pendubot}
Pendubot（Pendulum Robot）同为两连杆 RR 机器人，但\emph{仅第一关节有驱动}。它可看作小车倒摆的旋转版：主动的第一连杆扮演“小车”，去平衡被动的第二连杆（“摆”）。左端项与 Acrobot 相同，只是 $\mathbf{u}$ 作用在第一关节，是\emph{上驱动}形式：
\begin{equation}\label{eq:ua-pendubot}
  m_{11}\ddot{q}_1+m_{12}\ddot{q}_2+c_1+g_1=u,\qquad m_{21}\ddot{q}_1+m_{22}\ddot{q}_2+c_2+g_2=0.
\end{equation}
Acrobot 与 Pendubot 是同一机构“换驱动位置”的一对孪生例子，对比它们的可控性极具启发。

\subsection{反应轮摆（reaction wheel pendulum）}\label{sec:ua-rwp}
一根摆，远端装一个可被电机驱动的转盘（飞轮）。驱动飞轮产生的\emph{反作用力矩}带动摆运动。它是本章最简单的例子：若把 Acrobot 第二连杆配平使质心落在第二关节轴上（$\ell_{c2}=0$），Acrobot 即退化为反应轮摆。最简模型（把第一连杆质量集中于关节 2、并以相对水平度量飞轮角）为
\begin{equation}\label{eq:ua-rwp}
  J_1\ddot{q}_1+mg\ell\cos q_1=-u,\qquad J_2\ddot{q}_2=u,
\end{equation}
其中 $J_1=m_1\ell_{c1}^2+m_2\ell_1^2$、$J_2=I_2$、$m=m_1+m_2$。把两式相加消去 $u$ 又得等价下驱动形式 $J_1\ddot{q}_1+J_2\ddot{q}_2+mg\ell\cos q_1=0,\ J_2\ddot{q}_2=u$。

\begin{insight}[反应轮摆 $=$ 单摆 $\oplus$ 双积分器（并联）]
\cref{eq:ua-rwp} 把反应轮摆显示成一个单摆 $J_1\ddot{q}_1+mg\ell\cos q_1=-u$ 与一个双积分器 $J_2\ddot{q}_2=u$ 的\textbf{并联}：二者共享\emph{同一个}输入 $u$，符号相反。这个“两个无源子系统共享输入”的结构，是后面用无源性方法做 swing-up 的关键——既然并联的无源系统仍无源，只需为整体配一个合适的输出即可。
\end{insight}

\begin{note}[再点两个常被引用的例子]
\textbf{PVTOL}（平面垂直起降飞行器）有 3 自由度（平面位置 $+$ 姿态）、2 个推力输入，欠驱动度 1，是研究非最小相位飞行控制的标准模型。\textbf{TORA}（带旋转作动的平移振子）以一个旋转偏心质量去镇定一个弹簧滑块，常作无源性/反步法（\cref{ch:lyapunov-stability}）的试金石。本章不展开其推导，但其分析套路与上述例子同源，读者可按相同方法处理。
\end{note}

\section{平衡点与线性可控性}\label{sec:ua-controllability}

把 \cref{eq:ua-dyn} 写成状态空间形式。取状态 $\mathbf{x}=(\mathbf{x}_1,\mathbf{x}_2)=(\mathbf{q},\dot{\mathbf{q}})\in\mathbb{R}^{N},\ N=2n$，则
\begin{equation}\label{eq:ua-statespace}
  \dot{\mathbf{x}}=\boldsymbol{\mathcal{F}}(\mathbf{x},\mathbf{u})=\mathbf{f}(\mathbf{x})+\mathbf{G}(\mathbf{x})\mathbf{u},\quad
  \mathbf{f}=\begin{bmatrix}\mathbf{x}_2\\-\mathbf{M}^{-1}(\mathbf{C}\mathbf{x}_2+\boldsymbol{\phi})\end{bmatrix},\
  \mathbf{G}=\begin{bmatrix}\mathbf{0}\\\mathbf{M}^{-1}\mathbf{B}\end{bmatrix}.
\end{equation}

\begin{definition}[平衡点]\label{def:ua-equilibrium}
系统 $\dot{\mathbf{x}}=\boldsymbol{\mathcal{F}}(\mathbf{x},\mathbf{u})$ 的平衡点是满足 $\boldsymbol{\mathcal{F}}(\mathbf{x}_e,\mathbf{u}_e)=\mathbf{0}$ 的常向量 $(\mathbf{x}_e,\mathbf{u}_e)$。
\end{definition}

由 \cref{eq:ua-statespace}，平衡处必有 $\mathbf{x}_2=\dot{\mathbf{q}}=\mathbf{0}$，于是平衡构型由
\begin{equation}\label{eq:ua-eqcond}
  \boldsymbol{\phi}(\mathbf{q}_e)=\mathbf{B}\mathbf{u}_e
\end{equation}
决定。特别地当 $\mathbf{u}_e=\mathbf{0}$ 时，平衡点是势能的局部极值（极小或极大），因为 $\boldsymbol{\phi}$ 是势能梯度。若系统无势力项（无重力、无弹性），则 $\boldsymbol{\phi}\equiv\mathbf{0}$，整个构型空间的每一点 $(\mathbf{q},\mathbf{0})$ 都是平衡点——这与可控性密切相关，下文将见。

\begin{example}[Acrobot/Pendubot 的孤立平衡点]\label{ex:ua-acro-eq}
取 $u=0$，重力下 Acrobot/Pendubot 的全部平衡点只有四个（两连杆各自竖直向上/向下的组合）：
\[
  (q_1,q_2)\in\Big\{(-\tfrac{\pi}{2},0),\ (-\tfrac{\pi}{2},\pi),\ (\tfrac{\pi}{2},\pi),\ (\tfrac{\pi}{2},0)\Big\}.
\]
其中 $(\tfrac{\pi}{2},0)$ 是两连杆都竖直向上的\emph{倒立}平衡（swing-up 的目标）。
\end{example}

\subsection{线性可控性判据}\label{sec:ua-lincontrol}

\begin{definition}[平衡点处的线性可控性]\label{def:ua-lincontrol}
设 $(\mathbf{x}_e,\mathbf{u}_e)$ 是 $\dot{\mathbf{x}}=\boldsymbol{\mathcal{F}}(\mathbf{x},\mathbf{u})$ 的平衡点，记其线性化为 $\dot{\tilde{\mathbf{x}}}=\mathbf{F}\tilde{\mathbf{x}}+\mathbf{G}\tilde{\mathbf{u}}$，$\mathbf{F}=\tfrac{\partial\boldsymbol{\mathcal{F}}}{\partial\mathbf{x}}(\mathbf{x}_e,\mathbf{u}_e)$，$\mathbf{G}=\tfrac{\partial\boldsymbol{\mathcal{F}}}{\partial\mathbf{u}}(\mathbf{x}_e,\mathbf{u}_e)$。若线性系统 $(\mathbf{F},\mathbf{G})$ 可控，即
\begin{equation}\label{eq:ua-kalman}
  \operatorname{rank}\big[\,\mathbf{G},\ \mathbf{F}\mathbf{G},\ \mathbf{F}^2\mathbf{G},\ \dots,\ \mathbf{F}^{N-1}\mathbf{G}\,\big]=N,
\end{equation}
则称非线性系统在该平衡点\textbf{线性可控}。
\end{definition}

线性可控性允许在平衡点附近用线性反馈做局部指数镇定（如 LQR，\cref{sec:ctrl-lqr}），并为切换控制提供“近处控制器”。反之，在平衡点\emph{线性不可控}的系统（如下一章的非完整系统）即便做局部镇定也需根本不同的方法。

对机械系统 \cref{eq:ua-dyn}，平衡处 $\mathbf{x}_{1e}=\mathbf{q}_e,\ \mathbf{x}_{2e}=\mathbf{0}$，线性化矩阵为
\begin{equation}\label{eq:ua-linFG}
  \mathbf{F}=\begin{bmatrix}\mathbf{0}&\mathbf{I}\\-\mathbf{M}^{-1}(\mathbf{q}_e)\tfrac{\partial\boldsymbol{\phi}}{\partial\mathbf{q}}(\mathbf{q}_e)&\mathbf{0}\end{bmatrix},\qquad
  \mathbf{G}=\begin{bmatrix}\mathbf{0}\\\mathbf{M}^{-1}(\mathbf{q}_e)\mathbf{B}\end{bmatrix}.
\end{equation}

\begin{derivation}[为什么科氏项和 $\mathbf{M}^{-1}$ 的导数不出现]
对 \cref{eq:ua-statespace} 求 $\partial/\partial\mathbf{x}$。科氏/离心项 $\mathbf{C}(\mathbf{q},\dot{\mathbf{q}})\dot{\mathbf{q}}$ 关于速度是\emph{二次}的，其偏导在 $\dot{\mathbf{q}}=\mathbf{0}$ 处为零；而 $\mathbf{M}^{-1}$ 的偏导总乘以因子 $(\boldsymbol{\phi}-\mathbf{B}\mathbf{u})$，后者在平衡处（由 \cref{eq:ua-eqcond}）为零。故 $\mathbf{F}$ 只剩势力项梯度通过 $-\mathbf{M}^{-1}\partial\boldsymbol{\phi}/\partial\mathbf{q}$ 进入左下块。
\end{derivation}

把线性化按被动/主动分块（下驱动），得 $\mathbf{M}_{11}\ddot{\tilde{\mathbf{q}}}_1+\mathbf{M}_{12}\ddot{\tilde{\mathbf{q}}}_2=-\tfrac{\partial\boldsymbol{\phi}_1}{\partial\mathbf{q}}(\mathbf{q}_e)\tilde{\mathbf{q}}_1$ 等。由于 $\mathbf{M}(\mathbf{q}_e)$ 满秩，可控性要求 $\tfrac{\partial\boldsymbol{\phi}_1}{\partial\mathbf{q}}(\mathbf{q}_e)$（尺寸 $\ell\times n$，但仅 $\ell\times\ell$ 的被动块起约束作用）\emph{行满秩}。这给出一个干净的必要条件。

\begin{theorem}[线性可控性的必要条件]\label{thm:ua-lincontrol-nec}
下驱动系统在平衡点 $\mathbf{q}=\mathbf{q}_e,\dot{\mathbf{q}}=\mathbf{0},\mathbf{u}=\mathbf{u}_e$ 线性可控，\emph{仅当}
\begin{enumerate}
  \item $m\ge\ell$，即主动关节数不少于被动关节数；且
  \item $\tfrac{\partial\boldsymbol{\phi}_1}{\partial\mathbf{q}}(\mathbf{q}_e)$ 行满秩。
\end{enumerate}
对上驱动系统，把条件 2 换成 $\tfrac{\partial\boldsymbol{\phi}_2}{\partial\mathbf{q}}(\mathbf{q}_e)$ 行满秩。
\end{theorem}

\begin{insight}[物理含义：被动关节必须“吃得到”势力]
\cref{thm:ua-lincontrol-nec} 的条件 2 翻译成大白话：每个\emph{被动}关节轴在该平衡构型处必须有非零的势力（重力矩或弹性力矩）。直觉是——主动关节只能通过惯量耦合“隔山打牛”地影响被动关节的加速度；要在平衡点附近稳住被动关节，还得有一个随构型变化的恢复力把扰动“顶回来”，否则线性化后该方向毫无刚度。推论：无重力、无弹性的串联欠驱动机器人\emph{永远}不是线性可控的（如太空中的自由臂）。
\end{insight}

\begin{example}[反应轮摆：仅当 $mg\ell\ne0$ 可控]\label{ex:ua-rwp-control}
反应轮摆等价下驱动形式中 $\boldsymbol{\phi}_1=mg\ell\cos q_1$，故 $\partial\phi_1/\partial q_1=-mg\ell\sin q_1$。在倒立点 $q_{1e}=\pm\pi/2$ 处 $|\sin q_{1e}|=1$，条件 2 化为 $mg\ell\ne0$。以 $\mathbf{x}=(q_1,q_2,\dot{q}_1,\dot{q}_2)^\top$、归一化 $m=\ell=J_1=J_2=1,\ g=9.8$，倒立点线性化为 $\dot{\mathbf{x}}=\mathbf{F}\mathbf{x}+\mathbf{G}u$，
\[
  \mathbf{F}=\begin{bmatrix}0&0&1&0\\0&0&0&1\\9.8&0&0&0\\0&0&0&0\end{bmatrix},\quad
  \mathbf{G}=\begin{bmatrix}0\\0\\-1\\1\end{bmatrix}.
\]
该对可控（$mg\ell\ne0$ 时秩判据满足）。用 LQR（$\mathbf{Q}=\mathbf{I}_4,\ r=1$）解得镇定增益 $\mathbf{k}^\top=[-117.38,\,-30.38,\,-43.58,\,-13.43]$，$u=-\mathbf{k}^\top\mathbf{x}$ 即可在倒立点局部平衡。这正是后面切换控制里的“近处控制器”。
\end{example}

\begin{example}[Pendubot：水平时丢可控性]\label{ex:ua-pendubot-control}
Pendubot 是上驱动，被动的是第二连杆，$\phi_2=m_2\ell_{c2}g\cos(q_1+q_2)$，故条件要求 $\sin(q_{1e}+q_{2e})\ne0$。在平衡构型 $q_{1e}+q_{2e}=\pm\pi/2$（第二连杆竖直）处满足。但若进一步让第一连杆水平（$q_{1e}=0$），数值上控制性矩阵 \cref{eq:ua-kalman} 的行列式恰好过零——Pendubot 在“第一连杆水平、第二连杆竖直”这一特定构型处\emph{线性不可控}，在该区间其他平衡处均可控。这提醒我们：可控性是\emph{逐点}性质，必要条件满足不等于处处可控，须代入具体参数核验。
\end{example}

\section{部分反馈线性化（PFL）}\label{sec:ua-pfl}

既然欠驱动系统一般不能整体反馈线性化，退而求其次：用非线性反馈把\emph{能线性化的那部分}先化成双积分器。这就是部分反馈线性化。按“被线性化的是主动还是被动加速度”分为两类，二者都把系统化为含\emph{内动态}的二阶正规型，是后续设计的统一起点——可视作计算力矩法（\cref{sec:ctrl-robot-ctc}）在欠驱动情形的推广：全驱动时 $\ell=0$，PFL 退化为对全部坐标的反馈线性化。

\subsection{配置型 PFL（collocated）}\label{sec:ua-collocated}

“配置型”指在\emph{主动}关节加速度 $\ddot{\mathbf{q}}_2$ 与其输入 $\mathbf{u}$ 间建立线性关系。从下驱动的被动行 \cref{eq:ua-lower1} 出发：因 $\mathbf{M}_{11}\in\mathbb{R}^{\ell\times\ell}$ 由 $\mathbf{M}\succ0$ 知非奇异，解出被动加速度
\begin{equation}\label{eq:ua-q1ddot}
  \ddot{\mathbf{q}}_1=-\mathbf{M}_{11}^{-1}(\mathbf{M}_{12}\ddot{\mathbf{q}}_2+\mathbf{c}_1+\boldsymbol{\phi}_1),
\end{equation}
代入主动行 \cref{eq:ua-lower2}，得
\begin{equation}\label{eq:ua-coll-reduced}
  \bar{\mathbf{M}}_{22}\ddot{\mathbf{q}}_2+\bar{\mathbf{c}}_2+\bar{\boldsymbol{\phi}}_2=\mathbf{u},\qquad
  \begin{aligned}
    \bar{\mathbf{M}}_{22}&=\mathbf{M}_{22}-\mathbf{M}_{21}\mathbf{M}_{11}^{-1}\mathbf{M}_{12},\\
    \bar{\mathbf{c}}_2&=\mathbf{c}_2-\mathbf{M}_{21}\mathbf{M}_{11}^{-1}\mathbf{c}_1,\\
    \bar{\boldsymbol{\phi}}_2&=\boldsymbol{\phi}_2-\mathbf{M}_{21}\mathbf{M}_{11}^{-1}\boldsymbol{\phi}_1.
  \end{aligned}
\end{equation}

\begin{proposition}[$\bar{\mathbf{M}}_{22}$ 对称正定]\label{thm:ua-Mbar-spd}
$m\times m$ 矩阵 $\bar{\mathbf{M}}_{22}$ 在每个 $\mathbf{q}$ 处对称正定。
\end{proposition}
\begin{proof}
直接验证 $\bar{\mathbf{M}}_{22}=\mathbf{S}^\top\mathbf{M}\mathbf{S}$，其中
\begin{equation}\label{eq:ua-Smat}
  \mathbf{S}=\begin{bmatrix}-\mathbf{M}_{11}^{-1}\mathbf{M}_{12}\\\mathbf{I}_m\end{bmatrix}\in\mathbb{R}^{n\times m}
\end{equation}
对所有 $\mathbf{q}$ 秩为 $m$。$\mathbf{M}$ 对称正定且 $\mathbf{S}$ 列满秩，故 $\mathbf{S}^\top\mathbf{M}\mathbf{S}$ 对称正定。事实上 $\bar{\mathbf{M}}_{22}$ 正是 $\mathbf{M}_{22}$ 在 $\mathbf{M}$ 中的 \emph{Schur 补}（其正定性也可由 $\mathbf{M}\succ0$ 的 Schur 补判据直接得出）。
\end{proof}

由 \cref{thm:ua-Mbar-spd}，$\bar{\mathbf{M}}_{22}$ 可逆，故\textbf{配置型 PFL 控制律}
\begin{equation}\label{eq:ua-coll-law}
  \mathbf{u}=\bar{\mathbf{M}}_{22}\mathbf{a}_2+\bar{\mathbf{c}}_2+\bar{\boldsymbol{\phi}}_2
\end{equation}
（$\mathbf{a}_2\in\mathbb{R}^m$ 为外环控制）令 $\ddot{\mathbf{q}}_2=\mathbf{a}_2$。整个系统化为
\begin{equation}\label{eq:ua-coll-normalform}
  \boxed{\ \mathbf{M}_{11}\ddot{\mathbf{q}}_1+\mathbf{c}_1+\boldsymbol{\phi}_1=-\mathbf{M}_{12}\mathbf{a}_2,\qquad \ddot{\mathbf{q}}_2=\mathbf{a}_2.\ }
\end{equation}

\begin{definition}[二阶正规型]\label{def:ua-normalform}
\cref{eq:ua-coll-normalform} 称为\textbf{带输入驱动内动态的二阶正规型}，简称二阶正规型。第一式是\emph{内动态}（被动子系统），由外环 $\mathbf{a}_2$ 通过惯量耦合 $-\mathbf{M}_{12}\mathbf{a}_2$ 驱动；第二式是已被线性化的主动子系统。它与原系统反馈等价，是后续一切分析与设计的出发点。
\end{definition}

\begin{example}[小车倒摆的配置型 PFL]\label{ex:ua-cart-coll}
取下驱动形式（$q_1=\theta$ 被动、$q_2=x$ 主动）并归一化常数，动力学为
$\ddot{q}_1+\cos q_1\,\ddot{q}_2-\sin q_1=0$，$\cos q_1\,\ddot{q}_1+2\ddot{q}_2-\dot{q}_1^2\sin q_1=u$。配置型 PFL 控制
\[
  u=(2-\cos^2 q_1)\,a_2+\cos q_1\sin q_1-\dot{q}_1^2\sin q_1
\]
给出正规型 $\ddot{q}_1-\sin q_1=-\cos q_1\,a_2,\ \ddot{q}_2=a_2$。内动态是一个被 $a_2$ 经 $-\cos q_1$ 驱动的（倒）摆。
\end{example}

\subsection{非配置型 PFL（non-collocated）}\label{sec:ua-noncollocated}

更出人意料的是：在一定耦合条件下，可反过来线性化\emph{被动}关节加速度 $\ddot{\mathbf{q}}_1$。关键在于惯量阵的非对角块 $\mathbf{M}_{12}$（及 $\mathbf{M}_{21}=\mathbf{M}_{12}^\top$）产生的耦合广义力：主动输入不仅使主动关节加速，也通过 $\mathbf{M}_{12}$ 牵动被动关节。

\begin{definition}[强惯量耦合]\label{def:ua-strongcoupling}
设 $\mathcal{U}\subset\mathbb{R}^n$ 是构型空间的开集。系统 \cref{eq:ua-lower1,eq:ua-lower2} 在 $\mathcal{U}$ 上\textbf{强惯量耦合}，当且仅当
\begin{equation}\label{eq:ua-strongcoupling}
  \operatorname{rank}\mathbf{M}_{12}(\mathbf{q})=\ell\quad\forall\,\mathbf{q}\in\mathcal{U}.
\end{equation}
因 $\mathbf{M}_{12}$ 是 $\ell\times m$，该条件蕴含 $m\ge\ell$（主动关节数不少于被动关节数）。
\end{definition}

若 $\mathbf{M}_{12}$ 行满秩 $\ell$，则 $\mathbf{M}_{12}\mathbf{M}_{12}^\top\in\mathbb{R}^{\ell\times\ell}$ 可逆，定义右伪逆
\begin{equation}\label{eq:ua-pinv}
  \mathbf{M}_{12}^{\dagger}=\mathbf{M}_{12}^\top(\mathbf{M}_{12}\mathbf{M}_{12}^\top)^{-1}.
\end{equation}
由被动行 \cref{eq:ua-lower1} 解出 $\ddot{\mathbf{q}}_2=-\mathbf{M}_{12}^{\dagger}(\mathbf{M}_{11}\ddot{\mathbf{q}}_1+\mathbf{c}_1+\boldsymbol{\phi}_1)$，代入主动行得 $\bar{\mathbf{M}}_{21}\ddot{\mathbf{q}}_1+\bar{\mathbf{c}}_2+\bar{\boldsymbol{\phi}}_2=\mathbf{u}$，其中
\begin{equation}\label{eq:ua-noncoll-terms}
  \bar{\mathbf{M}}_{21}=\mathbf{M}_{21}-\mathbf{M}_{22}\mathbf{M}_{12}^{\dagger}\mathbf{M}_{11},\quad
  \bar{\mathbf{c}}_2=\mathbf{c}_2-\mathbf{M}_{22}\mathbf{M}_{12}^{\dagger}\mathbf{c}_1,\quad
  \bar{\boldsymbol{\phi}}_2=\boldsymbol{\phi}_2-\mathbf{M}_{22}\mathbf{M}_{12}^{\dagger}\boldsymbol{\phi}_1.
\end{equation}
$\bar{\mathbf{M}}_{21}$ 行满秩 $\ell$。于是\textbf{非配置型 PFL 控制律}
\begin{equation}\label{eq:ua-noncoll-law}
  \mathbf{u}=\bar{\mathbf{M}}_{21}\mathbf{a}_1+\bar{\mathbf{c}}_2+\bar{\boldsymbol{\phi}}_2
\end{equation}
给出非配置二阶正规型
\begin{equation}\label{eq:ua-noncoll-normalform}
  \boxed{\ \mathbf{M}_{12}\ddot{\mathbf{q}}_2+\mathbf{c}_1+\boldsymbol{\phi}_1=-\mathbf{M}_{11}\mathbf{a}_1,\qquad \ddot{\mathbf{q}}_1=\mathbf{a}_1.\ }
\end{equation}
此时\emph{被动}关节 $\mathbf{q}_1$ 被直接线性化为双积分器，内动态 \cref{eq:ua-noncoll-normalform} 第一式描述剩下的（主动）坐标。

\begin{example}[小车倒摆的非配置型 PFL]\label{ex:ua-cart-noncoll}
\cref{ex:ua-cart-coll} 的小车倒摆在 $-\tfrac{\pi}{2}<q_1<\tfrac{\pi}{2}$ 内满足强惯量耦合（$\cos q_1\ne0$）。非配置控制
\[
  u=2\tan q_1-\dot{q}_1^2\sin q_1-\frac{1+\sin^2 q_1}{\cos q_1}\,a_1
\]
给出 $\cos q_1\,\ddot{q}_2-\sin q_1=-a_1,\ \ddot{q}_1=a_1$，在该区间有效。注意控制律在 $q_1\to\pm\pi/2$ 处因 $\cos q_1\to0$ 而奇异——非配置 PFL 通常只\emph{局部}成立，配置型则\emph{全局}成立（$\bar{\mathbf{M}}_{22}$ 处处正定，无奇异），这是两者的重要区别。
\end{example}

\begin{pitfall}[概念误区：「PFL 把欠驱动变成了全驱动」]
PFL 只把\emph{一半}坐标的加速度化成了可指定的 $\mathbf{a}$，另一半仍由内动态自由演化、无法独立指令。系统的本质欠驱动性\emph{并未}消失，只是被搬进了 \cref{eq:ua-coll-normalform}/\cref{eq:ua-noncoll-normalform} 的内动态里。真正的控制难题——如何借内动态把状态导向目标——恰恰留在 $\mathbf{a}$ 对内动态的间接作用中，这才是后面能量法要解决的。
\end{pitfall}

\section{输出反馈线性化、零动态与最小相位}\label{sec:ua-zerodynamics}

PFL 是“以某个坐标为输出”的输出反馈线性化的特例。更一般地，取光滑输出 $\mathbf{y}=\mathbf{h}(\mathbf{q})\colon\mathbb{R}^n\to\mathbb{R}^p$，在二阶正规型 \cref{eq:ua-coll-normalform} 上对 $\mathbf{y}$ 求两次导：
\begin{equation}\label{eq:ua-yddot}
  \ddot{\mathbf{y}}=\underbrace{(\mathbf{J}_2-\mathbf{J}_1\mathbf{M}_{11}^{-1}\mathbf{M}_{12})}_{\bar{\mathbf{J}}}\,\mathbf{a}_2+\boldsymbol{\eta}(\mathbf{q},\dot{\mathbf{q}}),\quad \mathbf{J}_i:=\tfrac{\partial\mathbf{h}}{\partial\mathbf{q}_i},
\end{equation}
其中 $\bar{\mathbf{J}}$（$p\times m$）称\emph{解耦矩阵}，$\boldsymbol{\eta}=\dot{\mathbf{J}}_1\dot{\mathbf{q}}_1+\dot{\mathbf{J}}_2\dot{\mathbf{q}}_2-\mathbf{J}_1\mathbf{M}_{11}^{-1}(\mathbf{c}_1+\boldsymbol{\phi}_1)$。当 $\bar{\mathbf{J}}$ 行满秩（须 $m\ge p$）时，系统具\textbf{向量相对度二}：输出须微分两次输入才显现。用右伪逆 $\bar{\mathbf{J}}^{\dagger}$ 令 $\mathbf{a}_2=\bar{\mathbf{J}}^{\dagger}(\bar{\mathbf{a}}_2-\boldsymbol{\eta})$，即得解耦线性输出 $\ddot{\mathbf{y}}=\bar{\mathbf{a}}_2$，外环 $\bar{\mathbf{a}}_2$ 可任意镇定或跟踪。

\begin{definition}[零动态]\label{def:ua-zerodynamics}
对输出 $\mathbf{y}=\mathbf{h}(\mathbf{q})$，记 $\Gamma=\{(\mathbf{q},\dot{\mathbf{q}})\colon\mathbf{h}(\mathbf{q})=\mathbf{0},\ \dot{\mathbf{h}}=\tfrac{\partial\mathbf{h}}{\partial\mathbf{q}}\dot{\mathbf{q}}=\mathbf{0}\}$ 为\textbf{零动态流形}。把 $\mathbf{y}\equiv\mathbf{0}$ 强制施加后，系统限制在 $\Gamma$ 上的降阶动态称为\textbf{零动态}。镇定 $\mathbf{y}=\mathbf{0}$ 使 $\Gamma$ 成为受控不变流形。
\end{definition}

\subsection{零动态的计算与物理意义}\label{sec:ua-zd-compute}

取 $\mathbf{y}=\mathbf{q}_2$（$p=m$，解耦矩阵即 $\mathbf{I}_m$）。零动态由 $\mathbf{q}_2\equiv\mathbf{0},\dot{\mathbf{q}}_2\equiv\mathbf{0},\mathbf{a}_2\equiv\mathbf{0}$ 代入内动态 \cref{eq:ua-coll-normalform} 第一式得：
\begin{equation}\label{eq:ua-zd-lower}
  \tilde{\mathbf{M}}_{11}\ddot{\mathbf{q}}_1+\tilde{\mathbf{c}}_1+\tilde{\boldsymbol{\phi}}_1=\mathbf{0},\quad \tilde{(\cdot)}\text{ 表示在 }\mathbf{q}_2=\mathbf{0},\dot{\mathbf{q}}_2=\mathbf{0}\text{ 处取值}.
\end{equation}
这正是“把主动关节冻结在 $\mathbf{q}_2=\mathbf{0}$、只剩 $\ell$ 个被动关节”的降阶拉格朗日系统。其总能量 $E=\tfrac12\dot{\mathbf{q}}_1^\top\tilde{\mathbf{M}}_{11}\dot{\mathbf{q}}_1+\tilde{P}(\mathbf{q}_1)$ 沿零动态\emph{守恒}（$\dot{E}=0$）：零动态的轨迹就是降阶系统的等能量曲线。

\begin{example}[Acrobot 的零动态 $=$ 单摆]\label{ex:ua-acro-zd}
Acrobot 配置型正规型（$y=q_2$）的零动态，由 \cref{eq:ua-zd-lower} 与 \cref{eq:ua-acro-terms} 得
\[
  \tilde{m}_{11}\ddot{q}_1+n_1 g\cos q_1=0,\quad
  \tilde{m}_{11}=m_1\ell_{c1}^2+m_2(\ell_1^2+\ell_{c2}^2+2\ell_1\ell_{c2})+I_1+I_2,\ n_1=m_1\ell_{c1}+m_2\ell_1+m_2\ell_{c2},
\]
\emph{恰是一个简单单摆}。其相平面上几乎所有轨迹都是周期轨道，平衡解\emph{非}渐近稳定。
\end{example}

\begin{definition}[最小相位 vs.\ 非最小相位]\label{def:ua-minphase}
若系统对所选输出的零动态在原点\emph{渐近稳定}，称\textbf{最小相位}；否则称\textbf{非最小相位}。对最小相位系统，把输出镇定到零自动使整个状态收敛；对非最小相位系统则不然——内动态可能持续振荡甚至发散。
\end{definition}

\begin{insight}[非最小相位是欠驱动控制的“拦路虎”]
\cref{ex:ua-acro-zd} 的零动态是无阻尼单摆，其周期轨道意味着 Acrobot（以 $q_2$ 为输出）\textbf{非最小相位}：哪怕你把髋角 $q_2$ 完美镇定到 $0$，被动的第一连杆 $q_1$ 仍会绕着自然单摆轨道荡来荡去，停不下来。这就是为什么不能简单“令被动关节角误差趋零”了事，而必须诉诸能量整形——直接去塑造这条等能量轨道，把它推到通过倒立点的那一条上。
\end{insight}

\begin{note}[反应轮摆是可全反馈线性化的例外]
欠驱动系统一般不可全反馈线性化，但反应轮摆是例外。取输出 $y_1=J_1 q_1+J_2 q_2$（角动量型组合），逐次求导 $y_2=\dot{y}_1,\ y_3=\dot{y}_2=-\tfrac{mg\ell}{J_1}\cos q_1,\ y_4=\dot{y}_3=\tfrac{mg\ell}{J_1}\sin q_1\,\dot{q}_1$，到第四阶输入 $u$ 才显现。该输出在 $0<q_1<\pi$ 内\emph{相对度四 $=$ 系统阶数}，零动态为空，故可全反馈线性化为 Brunovsky 标准型 $\dot{y}_1=y_2,\dots,\dot{y}_4=a$。代价是控制律含因子 $1/\sin q_1$，在 $q_1\to0,\pi$ 处奇异，且要求初始时摆在水平线之上。
\end{note}

\begin{note}[虚拟完整约束（VHC）]
把输出 $\mathbf{h}(\mathbf{q})=\mathbf{0}$ 由\emph{反馈}（而非自然动力学或环境）维持，就得到一条\textbf{虚拟完整约束}：若初始在 $\Gamma$ 上，闭环轨迹永远留在 $\Gamma$ 上。VHC 是协调主动/被动关节做出指定运动的利器，在行走、蛇形、攀援、体操机器人中尤为有用。例如令 Acrobot 满足 $q_1+0.5q_2=0$ 可模拟连续接触的攀援（brachiation）运动——只要该 $\mathbf{h}$ 给出相对度二的输出，用输出反馈线性化即可施加。VHC 将在腿足运动控制中再次登场。
\end{note}

\section{能量整形与摆起（swing-up）}\label{sec:ua-energy}

现在解决核心问题：把摆从下垂\emph{泵}到倒立。回到总能量。对 \cref{eq:ua-dyn}，由拉格朗日系统的无源性，
\begin{equation}\label{eq:ua-passivity}
  \dot{E}=\dot{\mathbf{q}}^\top\mathbf{B}\mathbf{u},
\end{equation}
即系统是从输入 $\mathbf{B}\mathbf{u}$ 到输出 $\dot{\mathbf{q}}$ 的\emph{无源}映射。能量与无源性在此密不可分（全驱动鲁棒/自适应控制也建基于此）。能量法的精髓是：\emph{不去规划随时间的跟踪轨迹}，而是直接把能量这个标量调到目标值，让系统自己荡到对应的轨道上。

\subsection{单摆：能量泵的原型}\label{sec:ua-pendulum-pump}

考虑末端受力 $F$ 的单摆（可视作大系统中“被主动连杆甩动”的被动连杆），$\tau=\ell F$，
\begin{equation}\label{eq:ua-pend-eom}
  m\ell^2\ddot{\theta}+mg\ell\sin\theta=\tau,\qquad E=\tfrac12 m\ell^2\dot{\theta}^2+mg\ell(1-\cos\theta).
\end{equation}
$\tau=0$ 时能量守恒，等能集 $\Sigma_c=\{(\theta,\dot{\theta})\mid E=c\}$ 是不变流形（运动的第一积分）。相图中除两平衡点 $(0,0)$、$(\pm\pi,0)$ 外每条轨迹都周期；连接鞍点 $(\pm\pi,0)$（视为同一点）的那条特殊轨迹是\textbf{同宿轨道}。

\begin{definition}[同宿轨道]\label{def:ua-homoclinic}
动力系统的\textbf{同宿轨道}是连接一个鞍点平衡到其自身的轨迹，位于该鞍点稳定流形与不稳定流形的交集中。对单摆，它是恰好“能荡到倒立点（速度趋零）”的临界能量轨道，能量 $E_c=2mg\ell$（取下垂为零势能）。
\end{definition}

目标：设计 $\tau$ 使能量 $E(t)$ 收敛到给定 $c>0$。取 Lyapunov 候选
\begin{equation}\label{eq:ua-V-pend}
  V=\tfrac12(E-E_r)^2,\quad E_r=c,
\end{equation}
沿 \cref{eq:ua-pend-eom} 求导（用 $\dot{E}=\dot{\theta}\tau$）：
\begin{equation}\label{eq:ua-Vdot-pend}
  \dot{V}=(E-E_r)\dot{\theta}\,\tau.
\end{equation}
这说明系统从输入 $\tau$ 到输出 $y=(E-E_r)\dot{\theta}$ 无源。取
\begin{equation}\label{eq:ua-pump-law}
  \tau=-k\,y=-k(E-E_r)\dot{\theta},\quad k>0\ \Longrightarrow\ \dot{V}=-ky^2\le0.
\end{equation}

\begin{derivation}[LaSalle 给出收敛]
$\dot{V}\le0$ 仅半负定。$\dot{V}=0$ 当且仅当 $y=0$，即 $E=E_r$ 或 $\dot{\theta}=0$。在 $\{\dot{V}=0\}$ 上求最大不变集（LaSalle，\cref{ch:lyapunov-stability}）：可证所有轨迹收敛到\emph{能量为 $E_r$ 的轨道}或\emph{$\dot{\theta}=0$ 的孤立平衡}。$\dot{\theta}=0$ 不能排除——开环平衡 $(0,0),(\pm\pi,0)$ 因控制在 $\dot{\theta}=0$ 时为零而仍是闭环平衡；但下垂平衡 $(0,0)$ 在闭环下\emph{不稳定}（取 $E_r=E_c$ 时），故几乎所有轨迹被推向 $E=E_r$ 的轨道。取 $E_r=E_c=$ 同宿轨道能量，系统就被驱向那条恰好够到倒立点的轨道。
\end{derivation}

\begin{theorem}[能量整形 swing-up 的收敛性]\label{thm:ua-energy-conv}
对单摆 \cref{eq:ua-pend-eom} 施加能量注入律 \cref{eq:ua-pump-law}，则闭环全部轨迹收敛到集合 $\{E=E_r\}\cup\{\dot{\theta}=0\}$；除下垂/倒立的孤立开环平衡外，轨迹趋于等能量 $E_r$ 的周期轨道。取 $E_r$ 为同宿轨道能量 $E_c$，即把摆驱向通过倒立点的临界轨道。
\end{theorem}

\begin{insight}[能量泵的“正反馈”直觉]
看 \cref{eq:ua-pump-law}：当能量不足（$E<E_r$）时，$\tau=k(E_r-E)\dot{\theta}$ 与速度\emph{同向}——顺着运动方向推，做正功、灌能量；能量过量（$E>E_r$）时则\emph{逆向}阻尼、抽走能量。如此每个摆动周期净注入/抽出能量，把 $E$ 单调逼向 $E_r$。这与“跟踪某条预设轨迹”截然不同：我们只盯住一个标量 $E$，让物理自己挑轨道，因此天然鲁棒、无需轨迹规划。
\end{insight}

\subsection{饱和处理}\label{sec:ua-saturation}

无源法的一大优点是天然容纳输入饱和。设 $|\tau|\le M$，把 \cref{eq:ua-pump-law} 改为
\begin{equation}\label{eq:ua-sat-law}
  \tau=-\operatorname{sat}_M(k y)=-\operatorname{sat}_M\!\big(k(E-E_r)\dot{\theta}\big),\quad
  \operatorname{sat}_M(\sigma)=\begin{cases}M,&\sigma>M,\\\sigma,&|\sigma|\le M,\\-M,&\sigma<-M.\end{cases}
\end{equation}
饱和函数是“第一、三象限非线性”，即 $\sigma\,\operatorname{sat}_M(\sigma)\ge0$ 恒成立。故
\begin{equation}\label{eq:ua-Vdot-sat}
  \dot{V}=-y\,\operatorname{sat}_M(ky)\le0,
\end{equation}
LaSalle 结论不变。可见 swing-up 的能量律对饱和\emph{无须重新设计}，只要把增益夹进饱和即可——这是相较于反馈线性化类方法（饱和会破坏精确对消）的实在优势。

\subsection{反应轮摆的 swing-up}\label{sec:ua-rwp-swingup}

反应轮摆 \cref{eq:ua-rwp} 是“单摆 $\oplus$ 双积分器”的并联，两子系统各自从力矩到速度无源。记摆能量 $E_1=\tfrac12 J_1\dot{q}_1^2+mg\ell\sin q_1$、飞轮能量 $E_2=\tfrac12 J_2\dot{q}_2^2$，则 $\dot{E}_1=-\dot{q}_1 u,\ \dot{E}_2=\dot{q}_2 u$。取
\begin{equation}\label{eq:ua-rwp-V}
  V=\tfrac12(E_1-E_r)^2+E_2\ \Longrightarrow\ \dot{V}=\big(\dot{q}_2-(E_1-E_r)\dot{q}_1\big)u=:y\,u.
\end{equation}
系统从 $u$ 到 $y=\dot{q}_2-(E_1-E_r)\dot{q}_1$ 无源，取饱和控制
\begin{equation}\label{eq:ua-rwp-law}
  u=-\operatorname{sat}_M(ky)=-\operatorname{sat}_M\!\big(k(\dot{q}_2-(E_1-E_r)\dot{q}_1)\big),\quad k>0.
\end{equation}

\begin{proposition}[反应轮摆 swing-up 收敛]\label{thm:ua-rwp-conv}
取 $E_r>0$ 为摆能量参考。在 \cref{eq:ua-rwp} 上施加 \cref{eq:ua-rwp-law}，则闭环全部轨迹收敛到
\begin{equation}\label{eq:ua-rwp-set}
  \mathcal{M}=\{(\mathbf{q},\dot{\mathbf{q}})\mid (E_1-E_r)\cos q_1=0\}.
\end{equation}
即收敛到 $E_1=E_r$（此时由 \cref{eq:ua-rwp-law} 必有飞轮速度 $\dot{q}_2=0$）或 $\cos q_1=0$（$q_1=n\pi$、$\dot{q}_2=0$）。
\end{proposition}
\begin{proof}[证明思路]
$\dot{V}=-y\operatorname{sat}_M(ky)\le0$，对 $\{\dot{V}=0\}=\{y=0\}$ 用 LaSalle 求最大不变集即得 $\mathcal{M}$；细节为标准计算。取 $E_r$ 等于同宿轨道能量，即把摆驱向倒立点附近，随后切换到线性平衡控制器（如 \cref{ex:ua-rwp-control} 的 LQR）站稳。
\end{proof}

\subsection{Acrobot 的 swing-up 与平衡}\label{sec:ua-acro-swingup}

Acrobot 从配置型正规型 \cref{eq:ua-coll-normalform} 出发（$\ell=m=1$）：$m_{11}\ddot{q}_1+c_1+\phi_1=-m_{12}a_2,\ \ddot{q}_2=a_2$。关键在于内动态由外环 $a_2$ 驱动，故\emph{同一个} $a_2$ 既能镇定已线性化的主动子系统、又能塑造内动态能量。设计外环
\begin{equation}\label{eq:ua-acro-outer}
  a_2=-k_p q_2-k_d\dot{q}_2+k_3\operatorname{sat}_M\!\big((E-E_c)\dot{q}_1\big),
\end{equation}
其中 $E$ 是 Acrobot 总能量、$E_c$ 是倒立平衡处的能量。前两项 PD 把髋角 $q_2$ 拉向零（维持“身体不乱晃”的姿态约束），第三项是能量泵：把总能量推向 $E_c$，从而把第一连杆荡向倒立点。当状态进入倒立平衡的捕获域，即切换到 LQR 平衡控制（\cref{sec:ctrl-lqr}）站稳——典型仿真在约 $7.5\,$s 处完成切换。

\begin{pitfall}[易错点：能量到位 $\ne$ 状态到位]
能量律保证 $E\to E_c$，但 $E=E_c$ 是一张\emph{等能量面}，上面有无穷多状态，未必恰好是倒立点。能量法只能把状态送进“倒立点所在的那条临界轨道附近”，到了那里必须\emph{靠局部线性控制器抓住}。因此 swing-up 永远是“能量法 $+$ 切换 $+$ 局部镇定”的组合拳，缺一不可；单靠能量律无法稳定在倒立点。
\end{pitfall}

\section{切换控制：缝合全局甩起与局部平衡}\label{sec:ua-switching}

把上述思路总结成一个完整策略。倒立点附近系统线性可控（\cref{sec:ua-controllability}），可用 LQR 做局部指数镇定，但其有效范围（捕获域 / 吸引域）有限；远离倒立点时线性化失效，改用能量法把状态\emph{驱入}该捕获域。两者由一个切换逻辑缝合：

\begin{algo}[能量法 $+$ LQR 切换 swing-up]\label{algo:ua-switch}
\begin{enumerate}
  \item \textbf{建模与线性化}：把系统化为二阶正规型 \cref{eq:ua-coll-normalform}；在倒立平衡 $\mathbf{x}_e$ 处线性化得 $(\mathbf{F},\mathbf{G})$，核验线性可控（\cref{thm:ua-lincontrol-nec}）。
  \item \textbf{局部控制器}：对 $(\mathbf{F},\mathbf{G})$ 解 LQR（\cref{sec:ctrl-lqr}）得增益 $\mathbf{K}$，$\mathbf{u}_{\mathrm{bal}}=-\mathbf{K}(\mathbf{x}-\mathbf{x}_e)$。
  \item \textbf{全局甩起}：设计能量注入外环 \cref{eq:ua-acro-outer}（或对反应轮摆用 \cref{eq:ua-rwp-law}），参考能量 $E_r=E_c$ 取倒立点能量。
  \item \textbf{切换面}：定义捕获域 $\mathcal{B}=\{\mathbf{x}\colon (\mathbf{x}-\mathbf{x}_e)^\top\mathbf{P}(\mathbf{x}-\mathbf{x}_e)\le\rho\}$（$\mathbf{P}$ 为 LQR 的 Riccati 解，$\rho$ 由仿真/估计确定）。当 $\mathbf{x}\in\mathcal{B}$ 切换到 $\mathbf{u}_{\mathrm{bal}}$，否则用能量律。
\end{enumerate}
\end{algo}

\begin{insight}[捕获域：两套控制律的“接力区”]
切换控制的成败系于捕获域 $\mathcal{B}$ 的估计。$\mathcal{B}$ 太大，切进去时 LQR 的线性近似已失真、可能发散；太小，能量法的轨道又难以精确命中。实务上常用 Lyapunov 水平集（LQR 的 $\mathbf{x}^\top\mathbf{P}\mathbf{x}$ 取适当 $\rho$）作 $\mathcal{B}$，并辅以加滞回（hysteresis）防止在切换面上抖振。能量法的好处是它把状态\emph{反复}带到倒立点附近，给切换以多次机会，因此对 $\mathcal{B}$ 的精度不敏感、整体鲁棒。
\end{insight}

\begin{note}[能量整形 vs.\ 轨迹优化]
甩起也可用轨迹优化（如 DDP/iLQR，\cref{ch:ddp-ilqr} 已对“从静止甩起到直立”举例）：离线解一条最优轨迹再跟踪。两条路线互补——能量法\emph{在线、无需规划、对模型误差与饱和鲁棒}，但只控能量这一标量、对终端状态不精细；轨迹优化\emph{精细、可塞入约束与代价}，但依赖较准的模型且需在线/离线求解。工程上常以能量法做主甩起、优化法做精修，或干脆用本章的切换框架兜底。
\end{note}

\section{非线性可控性：被平衡点线性化“看漏”的可达性}\label{sec:ua-nonlinear-control}

最后澄清一个易混点：线性可控性（\cref{def:ua-lincontrol}）只看平衡点处线性近似的可控性，它可能\emph{严重低估}非线性系统的真实可达能力。

\begin{insight}[线性不可控 $\ne$ 不可达]
Acrobot 在倒立点线性可控（\cref{ex:ua-acro-eq} 附近满足 \cref{thm:ua-lincontrol-nec}），故能用 LQR 局部镇定——这是\emph{有}线性可控性的情形。但许多欠驱动系统在某些平衡点线性不可控（如 \cref{ex:ua-pendubot-control} 中第一连杆水平的 Pendubot，或下章的非完整轮式机器人），线性化丢掉了\emph{二阶及更高阶}的运动生成机制。然而它们在非线性意义下往往仍\textbf{（强）可达}：通过李括号 $[\mathbf{f},\mathbf{g}]$、$[\mathbf{g}_1,\mathbf{g}_2]$ 等“抖动”出新方向（如汽车反复打方向盘加进退实现横向平移）。判别这类可达性须用李括号/可达性分布（Chow 定理），属几何非线性控制范畴。
\end{insight}

\begin{note}[与全驱动可控性的对照，及承接下章]
全驱动机械臂在任意构型都（局部强）可控且可全反馈线性化，可控性几乎不成问题；欠驱动系统则把可控性推到舞台中央：要分清“线性可控（可用 LQR 局部镇定、可切换）”“线性不可控但非线性可达（须李括号/几何方法）”“真不可达”三层。本章处理的是有势力、可借能量整形与切换求解的一类；\emph{无漂移}（driftless）的非完整系统——线性不可控、靠李括号可达、且受 Brockett 定理所限\emph{无光滑时不变镇定律}——是另一类受限系统，留待 \cref{ch:nonholonomic} 专门展开。本章的能量/无源性视角还将与安全控制的控制 Lyapunov/障碍函数（\cref{ch:clf-cbf}）汇合：那里 swing-up 的 Lyapunov 候选自然升级为 CLF，捕获域升级为前向不变安全集。
\end{note}

\section{本章小结}\label{sec:ua-summary}

\begin{table}[H]
\centering
\caption{欠驱动系统与摆起控制：核心知识点速查}
\label{tab:ua-summary}
\small
\begin{tabular}{lll}
\toprule
知识点 & 核心条件/结论 & 关键应用 \\
\midrule
欠驱动度 & $\ell=n-m$；上/下驱动分块 & 柔性关节、腿足、空间机器人 \\
二阶约束 & 被动行右端为零 $\Rightarrow$ 不能跟踪任意轨迹 & 区分全驱动/欠驱动 \\
平衡点 & $\boldsymbol{\phi}(\mathbf{q}_e)=\mathbf{B}\mathbf{u}_e$；$\mathbf{u}_e{=}0$ 时为势能极值 & Acrobot 四平衡点 \\
线性可控（必要） & $m\ge\ell$ 且 $\partial\boldsymbol{\phi}_{1}/\partial\mathbf{q}$ 行满秩 & 被动关节须有非零势力 \\
配置型 PFL & $\mathbf{u}{=}\bar{\mathbf{M}}_{22}\mathbf{a}_2{+}\bar{\mathbf{c}}_2{+}\bar{\boldsymbol{\phi}}_2$；$\bar{\mathbf{M}}_{22}{=}\mathbf{S}^\top\mathbf{M}\mathbf{S}\succ0$ & 全局有效、$\ddot{\mathbf{q}}_2{=}\mathbf{a}_2$ \\
非配置型 PFL & 强惯量耦合 $\operatorname{rank}\mathbf{M}_{12}{=}\ell$；用 $\mathbf{M}_{12}^{\dagger}$ & 局部有效、$\ddot{\mathbf{q}}_1{=}\mathbf{a}_1$ \\
零动态 & 冻结输出后的降阶拉格朗日系统、$\dot{E}{=}0$ & Acrobot 零动态 $=$ 单摆 \\
非最小相位 & 零动态非渐近稳定（周期轨道） & 不能简单令输出趋零 \\
能量整形 & $V{=}\tfrac12(E{-}E_r)^2$，$\tau{=}{-}k(E{-}E_r)\dot{\theta}$ & swing-up，免轨迹规划 \\
饱和 & $\operatorname{sat}_M$ 第一三象限非线性、LaSalle 不变 & 实际力矩受限 \\
切换控制 & 远处能量法 $+$ 近处 LQR，捕获域 $\mathbf{x}^\top\mathbf{P}\mathbf{x}{\le}\rho$ & Acrobot/反应轮摆甩起平衡 \\
非线性可达 & 平衡点线性化丢可控；李括号/Chow 仍可达 & 对照非完整系统 \\
\bottomrule
\end{tabular}
\end{table}

主线一句话：\textbf{全驱动“抵消”动力学（计算力矩），欠驱动“利用”动力学（PFL $+$ 能量整形 $+$ 切换）}。PFL 把能线性化的坐标化为双积分器、把欠驱动性收进内动态；内动态的零动态常非最小相位，故须以能量/无源性直接塑造其轨道，把摆泵向同宿轨道；最后在倒立点切换到 LQR 站稳。

\section*{常见误解}
\addcontentsline{toc}{section}{常见误解}
\begin{enumerate}
  \item \textbf{“欠驱动只是执行器坏了”}：欠驱动常是\emph{有意}的设计或约束（接触、非完整、动量守恒）使然，并非故障；其控制理论自成体系。
  \item \textbf{“PFL 把系统变全驱动了”}：只线性化一半坐标，欠驱动性搬进内动态，本质难题不减。
  \item \textbf{“配置型与非配置型 PFL 对等”}：配置型 $\bar{\mathbf{M}}_{22}\succ0$ 处处可逆、\emph{全局}有效；非配置型依赖强惯量耦合、常\emph{局部}有效且在耦合退化处奇异。
  \item \textbf{“被动关节角误差趋零就稳了”}：若零动态非最小相位（如单摆），输出趋零时内动态仍周期振荡，状态不收敛。
  \item \textbf{“能量到 $E_c$ 就站住了”}：$E=E_c$ 是等能量面、含无穷多状态；能量法只把状态送到倒立点附近，须切换到局部镇定才稳定。
  \item \textbf{“线性不可控就没救”}：线性化会漏掉二阶可达性；许多线性不可控系统经李括号仍（强）可达，只是需几何方法（Chow 定理）。
  \item \textbf{“饱和会毁掉 swing-up”}：恰相反，能量律对饱和鲁棒——$\operatorname{sat}_M$ 的第一三象限性保住 $\dot{V}\le0$，无需重设计。
  \item \textbf{“切换面越大越好”}：捕获域太大会进入 LQR 线性近似失真区而发散；宜取 Lyapunov 水平集并加滞回防抖。
\end{enumerate}

\section*{延伸阅读}
\addcontentsline{toc}{section}{延伸阅读}
\begin{itemize}
  \item \textbf{Spong, Hutchinson \& Vidyasagar, \emph{Robot Modeling and Control}}\cite{spong2006robot}：本章主参考，欠驱动机器人一章，PFL/零动态/无源性 swing-up 的系统处理（第二版第 13 章内容更完整）。
  \item \textbf{Spong, ``The swing up control problem for the Acrobot''}\cite{spong1995swingup}：Acrobot 能量法甩起的原创论文。
  \item \textbf{Spong, ``Partial feedback linearization of underactuated mechanical systems''}\cite{spong1994pfl}：配置/非配置 PFL 与二阶正规型的来源。
  \item \textbf{Fantoni \& Lozano, \emph{Non-linear Control for Underactuated Mechanical Systems}}\cite{fantoni2002nonlinear}：欠驱动系统能量/无源性控制专著，含反应轮摆、Pendubot、TORA、PVTOL。
  \item \textbf{Block, Åström \& Spong, \emph{The Reaction Wheel Pendulum}}\cite{block2007reaction}：反应轮摆专论，建模到无源性 swing-up 一气呵成。
  \item \textbf{Westervelt et al., \emph{Feedback Control of Dynamic Bipedal Robot Locomotion}}\cite{westervelt2007feedback}：虚拟完整约束（VHC）与混杂零动态在双足行走中的奠基性应用。
\end{itemize}
